% ============================================================
%  Motivation / Introduction with \cite commands inserted.
%  Citation keys resolve against references.bib.
%  Only citations were added; wording is unchanged from your draft
%  except where noted in a % comment.
% ============================================================

\section{Motivation}

Modular self-reconfigurable robots derive their utility from changing topology
at runtime: modules dock into larger assemblies, undock, and rearrange into new
morphologies for new tasks~\cite{yim2007modular,thalamy2019survey}. Simulation is
the primary means by which such systems are designed, evaluated, and controlled
prior to hardware. The difficulty is not the dynamics of any single
configuration. It is that configuration is itself a runtime variable, and
representing that change faithfully couples the physics engine, the robot model,
and the control logic.

High-fidelity engines such as MuJoCo~\cite{todorov2012mujoco}, Isaac
Sim~\cite{nvidia_isaacsim}, and Gazebo~\cite{koenig2004gazebo} already integrate
multi-body robots with mature contact and actuation models. A \emph{fixed}
assembly is therefore comparatively well served. A dock or undock, however,
changes the identity of the mechanism during the run: independent bodies become
one, degrees of freedom appear or vanish, and the logical claim that two modules
now constitute a single assembly must remain consistent with the constraint the
engine actually enforces. General-purpose engines do not represent that second
half. Four difficulties follow.

First, topology is dynamic, yet most high-fidelity engines fix it at compile
time. MuJoCo~\cite{todorov2012mujoco} compiles bodies, joints, and constraint
slots once; model topology is immutable thereafter, so a dock cannot create a
constraint. The practical alternative is a pre-reserved pool of inactive welds
that are claimed, re-pointed, and released. Adding a joint on contact is not an
engine operation. Recompilation is expensive and breaks state continuity;
silently fabricating a connection is incorrect. Correct pool accounting,
including refusal when the pool is exhausted, is therefore unavoidable systems
work.

Second, every dock is simultaneously a physical event and a semantic event, and
the two must agree. A logical connection should be committed only after the
backend confirms that the corresponding constraint exists. Otherwise the
connectivity graph can assert a latch that physics never formed, or the converse,
and every downstream controller inherits the inconsistency. Scripts that set a
Boolean and add a joint in the same step provide neither confirmation nor
atomicity, nor a defined failure mode. Intra-step order is equally consequential:
releases must precede new detection so that a connection cannot break and reform
within a single step.

Third, modular semantics are expressed between connector frames, whereas physics
constraints are expressed between body frames. A dock specifies a relative pose
\(T_{\mathrm{connA}\leftarrow\mathrm{connB}}\); a weld constrains
\[
T_{\mathrm{bodyA}\leftarrow\mathrm{bodyB}}
= T_{\mathrm{bodyA}\leftarrow\mathrm{connA}}\,
T_{\mathrm{connA}\leftarrow\mathrm{connB}}\,
T_{\mathrm{connB}\leftarrow\mathrm{bodyB}}.
\]
The conversion is stricter when a connector lies on an articulated child rather
than the module root: the child's pose relative to the root must be measured and
mapped back to a root free-joint pose before the weld is activated. Treating the
child as the root, or composing frames in the wrong order, injects a large pose
error that is easily misread as numerical instability rather than incorrect
geometry.

Fourth, correctness cannot be read from the physics alone. Compatibility,
acceptance regions, assembly membership, mating orientation, and the event log
are semantic judgments. An engine will weld incompatible connectors at an
arbitrary pose; only a layer above the integrator can reject that outcome. The
mating orientation in particular must be recorded, because configuration
recognition and reconfiguration planning depend on the discrete mating
state~\cite{liu2019distributed}. Verification further requires a notion of ``the
same decision'' that is independent of exact poses: once dynamics act, poses may
diverge, but docking decisions must not. When semantics and physics are fused in
a single script, there is no independent oracle against which a frame-convention
error can be distinguished from a genuine dynamical difference.
% NOTE: the \cite{liu2019distributed} sentence above is an optional addition;
% delete it if you want your paragraph exactly as drafted.

Existing tools address only one side of this split. Lattice and behavioral
simulators such as VisibleSim~\cite{thalamy2022visiblesim} treat reconfiguration
semantics as first-class and scale to very large ensembles, but they omit
rigid-body dynamics by design. Dedicated physics-based modular simulators such as
ReMod3D~\cite{collins2013remod3d} combine docking with dynamics, yet they are
bound to a single engine and offer no portable backend contract. Frameworks such
as Revolve2~\cite{hupkes2018revolve,revolve2software} model modular bodies on
MuJoCo but target the evaluation of fixed morphologies; runtime dock and undock
lie outside their scope. General engines and kinematics
libraries~\cite{pinciroli2012argos,iory_scikitrobot,coumans2021pybullet} provide
dynamics without a native notion of module, connector, or dock, so each of the
four difficulties is left as engine-specific glue. No existing system jointly
owns faithful physics and portable, verifiable reconfiguration semantics.

The problem is therefore to keep a changing logical topology in exact, checkable
agreement with a fixed-topology physics engine across the instant of change,
without binding modular meaning to one engine. This paper argues that the
appropriate response is a semantic layer that owns topology, docking acceptance,
and the dock/undock lifecycle, and that delegates physics to a swappable backend.
The remainder of the paper presents that layer: a portable module description, a
narrow adapter contract, a weld-pool bridge onto compile-time-fixed engines, and
a differential conformance suite that runs identical scenarios on a
dependency-free kinematic backend and on the physics backend and asserts they
agree on connector frames, docking verdicts, and the committed event sequence,
so that an adapter error surfaces as a semantic disagreement rather than a
subtle dynamical difference.
